\chapter{CONCLUSÃO}

Este trabalho apresentou o desenvolvimento da plataforma NaraPsi, destinada à organização de registros e documentos relacionados ao atendimento psicológico. A solução reuniu autenticação, controle de acesso, persistência de dados, armazenamento de arquivos, auditoria, integração com assinatura eletrônica e apoio à produção textual por modelo de linguagem. A arquitetura e os controles implementados foram descritos em relação aos princípios de segurança e proteção de dados, sem que essa descrição constitua certificação de conformidade integral com a LGPD.

O módulo de inteligência artificial utilizou o Qwen2.5-3B-Instruct quantizado, executado localmente por meio de \texttt{llama-cpp-python}. Essa configuração permitiu processar o contexto no ambiente da aplicação, sem encaminhá-lo a um provedor externo de inferência. No protocolo exploratório, cinco dos seis casos atenderam integralmente aos critérios definidos. A latência média foi de 12,750~s, no equipamento utilizado, e uma das seis respostas apresentou informações não sustentadas pelo contexto. Na segmentação manual, duas de dez proposições clínicas verificáveis foram classificadas como sem suporte.

Esses resultados demonstram a viabilidade técnica da geração local, mas não autorizam concluir que as respostas sejam isentas de alucinações ou adequadas ao uso clínico autônomo. A amostra foi reduzida, cada caso de IA foi executado uma única vez e a classificação foi realizada pelo próprio autor. Os oito casos funcionais e de integração atenderam aos critérios definidos no ambiente de desenvolvimento; contudo, a execução manual e o conjunto restrito de cenários não substituem uma suíte automatizada nem uma validação independente em ambiente de produção.

Conclui-se que a NaraPsi constitui um protótipo capaz de integrar gestão documental e apoio textual local em um mesmo fluxo, desde que a saída da inteligência artificial seja tratada como minuta e submetida à revisão do profissional responsável. Como trabalhos futuros, recomenda-se automatizar os testes funcionais, ampliar e repetir o protocolo de IA, empregar avaliadores independentes, comparar configurações de inferência e validar a plataforma com usuários em ambiente controlado. Também devem ser aprofundadas as medidas organizacionais necessárias à proteção de dados, incluindo gestão de retenção, direitos dos titulares e resposta a incidentes.
